\documentclass[10pt,a4paper]{amsart}
\usepackage[margin=19mm]{geometry}
\usepackage{amsmath,amssymb,mathtools}
\usepackage[scr=rsfs,cal=euler]{mathalfa}
\usepackage{xfrac}
\usepackage[unicode,hidelinks,bookmarks=false]{hyperref}
\newcommand{\ie}{i.e.\ }
\newcommand{\cf}{cf.\ }
\newcommand{\resp}{respectively\ }
\newcommand{\Subsection}{\subsection}
\providecommand{\nobreakdash}{\nobreak\hbox{-}}
\setlength{\emergencystretch}{2em}
\multlinegap0pt
\usepackage{amssymb}
\usepackage{tikz}
\usepackage[shortlabels]{enumitem}
\newtheorem{lemma}{Lemma}
\newtheorem{theorem}{Theorem}
\newcommand\Lapd{\Delta_{d_1, d_2}}
\newcommand\Lapdd{\Delta_{d}}
\title{Herz versus Fefferman: Symmetric and asymmetric Bochner--Riesz theory}
\author{Peng Chen, Dangyang He, Adam Sikora, Lixin Yan}
\date{Source: arXiv:2608.09247v1 (CC BY 4.0)}
\begin{document}
\maketitle

\noindent\emph{Source and licence.} Herz versus Fefferman: Symmetric and asymmetric Bochner--Riesz theory. Version arXiv:2608.09247v1 (CC BY 4.0), distributed by arXiv under the Creative Commons Attribution 4.0 International licence, http://creativecommons.org/licenses/by/4.0/, as recorded in the arXiv licence metadata for this exact version and shown on the abstract page https://arxiv.org/abs/2608.09247v1. Full licence notice: \texttt{licenses/arxiv-2608.09247-CC-BY-4.0.txt}.\par\smallskip

\noindent\textit{Editorial scope.} Counted unit: the article's theorem fefferman\_nondoubling (source0.tex lines 2780--2786). The article's complete original proof of it is delivered with it (source0.tex lines 2789--3017, one \texttt{proof} environment of 229 lines). No mathematics is written, repaired or re-derived: the statement and the proof are the article's own lines, in the article's own notation, and no retained line is substituted or re-flowed.  Retained uncounted as the source-local context the proof needs: lines 1177--1198; lines 1387--1397; lines 1696--1702.  Nothing was omitted from any retained span, so the package carries no omission record.  No retained line contains a citation, so the wrapper carries no bibliography and no bibliography entry.  No retained line contains a figure environment, an image-inclusion command or drawing code, so no image is delivered and the package declares no asset dependency.  Editorial formatting only, in addition: an AMS \texttt{amsart} preamble that restates the article's own declarations and macros used by the retained lines, namely the environments \texttt{lemma}, \texttt{theorem} on separate counters, together with the standard packages those lines need.  The retained environments print in this document as Lemma 1, Theorem 1 in order of appearance, whereas the article prints the same items with the numbering of its own full document.  The complete native source of the article as distributed in the arXiv e-print for this exact version is bundled unchanged as \texttt{sources/207227/source0.tex}.
\begin{lemma}\label{Nix}
Let $n_1,n_2>1$, and let $K$ be the integral operator with kernel
\begin{equation}
    K(x,y)=
    \begin{cases}
       x^{-\alpha}y^{-\beta}, & 1\le x\le y,\\
       x^{-\alpha'}y^{-\beta'}, & x>y,
    \end{cases}
\end{equation}
where $\alpha+\beta=\alpha'+\beta'$. If
\( p(\alpha+\beta-n_1)\ge n_2-n_1
\)
and
\begin{equation*}
    \frac{n_2}{n_2\wedge \alpha'}
    < p <
    \frac{n_1}{0\vee(n_1-\beta)},
\end{equation*}
then the corresponding operator $K$ is bounded from
$L^p([1,\infty),r^{n_1-1}\,dr)$ to $
    L^p([1,\infty),r^{n_2-1}\,dr).$
\end{lemma}

\begin{theorem}\label{herz_kl}
Let $d_1,d_2\ge 1$ and set $D=d_1\vee d_2$. Then
\[
    \sup_{R>0}\left\|E_{\sqrt{\Lapd}}^{kl}(R)\right\|_{L^p\to L^p}
    <\infty
\]
for all
\[
    \left|\frac{1}{p}-\frac{1}{2}\right|<\frac{1}{2D}.
\]
\end{theorem}

\begin{theorem}\label{herz_doubling}
Let $d>2$. Then 
\begin{align*}
    \sup_{R>0} \| E_{\sqrt{\Lapdd}}(R)\|_{p \to p} \le C \quad \textit{if} \quad 
        \left| \frac{1}{2} - \frac{1}{p} \right| < \frac{1}{2d}.
\end{align*}
\end{theorem}

  \begin{theorem}\label{fefferman_nondoubling}
Let $d_1,d_2>2$ and $d_1 \ne d_2$. Then
\begin{align*}
    \sup_{R>0} \left\| E_{\sqrt{\Lapd}}(R) \right\|_{p \to p} \le C \quad \textit{iff} \quad p=2.
\end{align*}

\end{theorem}

\begin{proof}
Thanks to Theorem~\ref{herz_kl}, we again only need to handle the $kk$ part of the kernel. To do so, we further decompose the $kk$ part into the diagonal part (up to complex constants)
\begin{align*}
    E_{\mathrm{on}}(R)(x,y) = \sum_{\ell=1}^2 \sum_{\pm} E_{\mathrm{on}}^{\ell,\pm}(R)(x,y),
\end{align*}
and the off-diagonal part
\begin{align*}
    E_{\mathrm{off}}(R)(x,y) = \sum_{\ell \ne j} \sum_{\pm} E_{\mathrm{off}}^{\ell j, \pm}(R)(x,y).
\end{align*}
The key point is the off-diagonal part. Indeed, we note that 
\begin{itemize}
    \item first, for $\ell=1,2$, $E_{\mathrm{on}}^{\ell,\pm}(R)$ acts as an operator from $L^p([1,\infty),|r|^{d_\ell-1} dr)$ to $L^p([1,\infty),|r|^{d_\ell-1} dr)$;
    \item second, the $kk$ term in its kernel takes the form $k_\ell(\pm i\lambda|x|)k_\ell(\pm i\lambda|y|)$;
    \item third, although the definition of the coefficient $B_1$ (resp. $B_2$) depends on $d_1$ and $d_2$, its asymptotic behaviour depends on $d_1$ (resp. $d_2$) only. 
\end{itemize}
The above three observations allow us to apply Theorem~\ref{herz_doubling} to the on-diagonal part $E_{\mathrm{on}}(R)$, yielding the following lemma.



\begin{lemma}\label{fefferman-on-diag}
Let \(D=d_1\vee d_2\). If
\(\left|1/p-1/2\right|<\frac1{2D}\), then
\[
\sup_{R>0}\|E_{\mathrm{on}}(R)\|_{L^p\to L^p}\leq C.
\]
\end{lemma}

Therefore, to show the unboundedness of $E_{\sqrt{\Lapd}}(R)$ on $L^p$ ($p\ne 2$), one only needs to prove
\begin{equation}\label{claim}
    \| E_{\mathrm{off}}(1)\|_{p \to p} = \infty,\quad \forall p\ne 2.
\end{equation}
In what follows,
we construct a counterexample establishing this assertion.

Note that the kernel of $E_{\mathrm{off}}(1)$ has the form (up to a constant)
\begin{align}\label{off-diag}
    \sum_{\ell \ne j} \left[ E_{\mathrm{off}}^{\ell j,+}(1)(x,y) - E_{\mathrm{off}}^{\ell j,-}(1)(x,y)  \right].
\end{align}
Suppose $|x|\ge |y|$, so that $0<|x|^{-1}\le |y|^{-1} \le 1$. We split the integration in $\lambda$ from \eqref{off-diag} into $\int_0^1 = \int_0^{|x|^{-1}} + \int_{|x|^{-1}}^{|y|^{-1}} + \int_{|y|^{-1}}^1$ (and $\int_0^1 = \int_0^{|y|^{-1}} + \int_{|y|^{-1}}^{|x|^{-1}} + \int_{|x|^{-1}}^1$ if $|x|\le |y|$). Define
\begin{align*}
    I(x,y) = \begin{cases}
        \int_0^{|x|^{-1}} \dots, & |x|\ge |y|,\\
        \int_0^{|y|^{-1}}\dots, & |x|\le |y|,
    \end{cases}\quad
    II(x,y) = \begin{cases}
        \int_{|x|^{-1}}^{|y|^{-1}} \dots, & |x|\ge |y|,\\
    \int_{|y|^{-1}}^{|x|^{-1}}\dots, & |x|\le |y|,
    \end{cases}\quad
    III(x,y) = \begin{cases}
        \int_{|y|^{-1}}^1 \dots, & |x|\ge |y|,\\
    \int_{|x|^{-1}}^1\dots, & |x|\le |y|.
    \end{cases}
\end{align*}
We then claim that $I+II$ acts as a bounded operator on $L^p$ for $p\in \left( \frac{2D}{D+1}, \frac{2D}{D-1} \right)$. The proof is similar to the argument from Theorem~\ref{herz_doubling}. By asymptotic formulas, one has (for $|x|\ge |y|$)
\begin{align*}
     \int_0^{|x|^{-1}} \lambda A(\pm i\lambda) k_\ell(\pm i\lambda|x|) k_j(\pm i\lambda|y|) d\lambda \sim \int_0^{|x|^{-1}} |x|^{2-d_\ell} |y|^{2-d_j} \lambda^{d_1+d_2+1-d_\ell-d_j} d\lambda \simeq |x|^{-d_\ell} |y|^{2-d_j}.
\end{align*}
Thus, $|I| = O(|x|^{-d_\ell} |y|^{2-d_j})$. As for $II$, 
\begin{align*}
    \int_{|x|^{-1}}^{|y|^{-1}} \lambda A(\pm i\lambda) k_\ell(\pm i\lambda|x|) k_j(\pm i\lambda|y|) d\lambda \sim |x|^{-d_\ell} |y|^{2-d_j} \int_1^{|x|/|y|} s^{\frac{d_\ell-1}{2}} e^{\mp i s} ds.
\end{align*}
Integration by parts yields
\begin{equation*}
    \left|\int_1^{|x|/|y|} s^{\frac{d_\ell-1}{2}} e^{\mp i s} ds\right|\lesssim \left( \frac{|x|}{|y|} \right)^{\frac{d_\ell-1}{2}}.
\end{equation*}
Hence, it follows by symmetry that
\begin{align*}
    |(I+II)(x,y)|\lesssim \begin{cases}
        |x|^{-\left( d_\ell+\frac{d_j-1}{2}-2 \right)} |y|^{-\frac{d_j+1}{2}}, & |y|\ge |x|,\\
        |x|^{-\frac{d_\ell+1}{2}} |y|^{-\left( d_j + \frac{d_\ell-1}{2} -2 \right)}, & |y|\le |x|.
    \end{cases}
\end{align*}
By Lemma~\ref{Nix}, $I+II$ is bounded from 
\begin{equation*}
    L^p([1,\infty),r^{d_j-1}dr) \to L^p([1,\infty),r^{d_\ell-1}dr)
\end{equation*}
for all
\begin{equation}
    \frac{d_\ell}{\min \left( d_\ell, \frac{d_\ell+1}{2} \right)} < p < \frac{d_j}{\max\left( 0, d_j - \frac{d_j+1}{2} \right)} \iff \frac{2d_\ell}{d_\ell+1} < p < \frac{2d_j}{d_j-1}.
\end{equation}
Since the stated $D$-dependent range is contained in this interval for each $\ell\neq j$, the claim follows.

As a consequence, it suffices to establish the $L^p$-unboundedness of $III$. Note that for $|x|\ge |y|$,
\begin{align*}
    \int_{|y|^{-1}}^1  \lambda A(\pm i\lambda) k_\ell(\pm i\lambda|x|) k_j(\pm i\lambda|y|) d\lambda \sim |x|^{-\frac{d_\ell-1}{2}} |y|^{-\frac{d_j-1}{2}}\int_{|y|^{-1}}^1 \lambda^{\alpha} e^{\mp i t\lambda} d\lambda := |x|^{-\frac{d_\ell-1}{2}} |y|^{-\frac{d_j-1}{2}} \mathcal{I}_\alpha(x,y),
\end{align*}
where $\alpha = \frac{d_\ell+d_j}{2}-2>0$ and $t = |x|+|y|$. Integration by parts twice gives
\begin{align}\nonumber
    \mathcal{I}_\alpha(x,y) &= \frac{e^{\mp it}}{it} \pm \frac{|y|^{-\alpha} e^{\mp it/|y|}}{it} + \frac{\alpha e^{\mp it}}{t^2} - \frac{\alpha e^{\mp it/|y|} |y|^{1-\alpha}}{t^2}    - \frac{\alpha (\alpha-1)}{t^2} \int_{|y|^{-1}}^1 e^{\mp it \lambda}  \lambda^{\alpha-2} d\lambda \\ \label{I1}
    &= \frac{e^{\mp i(|x|+|y|)}}{i(|x|+|y|)} + O\left(\frac{1}{|y|^{\alpha}(|x|+|y|)}\right) + O\left(\frac{1}{(|x|+|y|)^2}\right) + O\left(\frac{1}{|y|^{\alpha-1}(|x|+|y|)^2}\right).
\end{align}
Symmetrically, for $|x|\le |y|$,
\begin{align}\label{I2}
    \mathcal{I}_\alpha(x,y)= \frac{e^{\mp i(|x|+|y|)}}{i(|x|+|y|)} + O\left(\frac{1}{|x|^{\alpha}(|x|+|y|)}\right) + O\left(\frac{1}{(|x|+|y|)^2}\right) + O\left(\frac{1}{|x|^{\alpha-1}(|x|+|y|)^2}\right).
\end{align}

\textit{Claim~1. The contributions of the last three terms from \eqref{I1} and \eqref{I2} in $III$ define bounded operators
$$L^p([1,\infty),r^{d_j-1}dr) \to L^p([1,\infty),r^{d_\ell-1}dr) $$
for $\frac{2D}{D+1}<p<\frac{2D}{D-1}$.}

To prove \textit{Claim 1}, one only needs to show the $L^p$-boundedness of $S_1, S_2, S_3$, where
\begin{align*}
    &S_1(x,y) = \begin{cases}
        |x|^{-\frac{d_\ell-1}{2}-\alpha}|y|^{-\frac{d_j-1}{2}}(|x|+|y|)^{-1}, & |y|\ge |x|,\\
        |x|^{-\frac{d_\ell-1}{2}} |y|^{-\frac{d_j-1}{2}-\alpha} (|x|+|y|)^{-1}, & |y|\le |x|,
    \end{cases}  \\
    &S_2(x,y) = |x|^{-\frac{d_\ell-1}{2}}|y|^{-\frac{d_j-1}{2}} (|x|+|y|)^{-2},\\
    &S_3(x,y) = \begin{cases}
        |x|^{-\frac{d_\ell-1}{2}+1-\alpha}  |y|^{-\frac{d_j-1}{2} }(|x|+|y|)^{-2}, & |y|\ge |x|,\\
         |x|^{-\frac{d_\ell-1}{2}} |y|^{-\frac{d_j-1}{2} +1-\alpha }(|x|+|y|)^{-2}, & |y|\le |x|.
    \end{cases}
\end{align*}

\begin{proof}[Proof of Claim~1]
One checks easily that 
\begin{align*}
&|S_1(x,y)| \lesssim \begin{cases}
        |x|^{-\frac{d_\ell-1}{2}-\alpha} |y|^{-\frac{d_j+1}{2}}, & |y| \ge |x|,\\
        |x|^{-\frac{d_\ell+1}{2}} |y|^{-\frac{d_j-1}{2}-\alpha}, & |y|\le |x|,
    \end{cases} \quad     |S_2(x,y)| \lesssim \begin{cases}
        |x|^{-\frac{d_\ell-1}{2}} |y|^{-\frac{d_j+3}{2}}, & |y| \ge |x|,\\
        |x|^{-\frac{d_\ell+3}{2}} |y|^{-\frac{d_j-1}{2}}, & |y|\le |x|,
    \end{cases}\\
    &|S_3(x,y)| \lesssim \begin{cases}
        |x|^{-\frac{d_\ell-1}{2}+1-\alpha} |y|^{-\frac{d_j+3}{2}}, & |y| \ge |x|,\\
        |x|^{-\frac{d_\ell+3}{2}} |y|^{-\frac{d_j-1}{2}+1-\alpha}, & |y|\le |x|.
    \end{cases}
\end{align*}
Note that the  operators $S_1,S_2,S_3$ are all of homogeneous type. Moreover, for all $p\ge 1$, we have
\begin{align*}
    p \left(\frac{d_\ell-1}{2} + \alpha + \frac{d_j+1}{2} - d_j \right) = p (d_\ell-2) \ge d_\ell-d_j.
\end{align*}
Therefore, by Lemma~\ref{Nix}, $S_1$ is bounded in $L^p$ for all
\begin{align}\label{S1}
    \frac{2d_\ell}{2d_\ell \wedge (d_\ell+1)} < p < \frac{2d_j}{0 \vee (d_j-1)} \iff \frac{2d_\ell}{d_\ell+1} < p < \frac{2d_j}{d_j-1}.
\end{align}
Next, to apply Lemma~\ref{Nix} to $S_2$, one needs
\begin{equation*}
    p\left( \frac{d_\ell-1}{2} + \frac{d_j+3}{2} - d_j \right) \ge d_\ell-d_j,
\end{equation*}
which is satisfied if
\begin{align*}
    \begin{cases}
        1\le p \le \infty, & d_\ell-d_j=-2, \\
        p\ge \frac{2(d_\ell-d_j)}{d_\ell-d_j+2}, & d_\ell-d_j>-2,\\
        p \le \frac{2(d_j-d_\ell)}{d_j-d_\ell-2}, & d_\ell-d_j<-2.
    \end{cases}
\end{align*}
A straightforward calculation shows that the above assumption holds if $\frac{2d_\ell}{d_\ell+1}<p<\frac{2d_j}{d_j-1}$. So, Lemma~\ref{Nix} guarantees that $S_2$ is $L^p$-bounded for
\begin{equation}\label{S2}
    \frac{2d_\ell}{2d_\ell \wedge (d_\ell+3)} <p< \frac{2d_j}{0 \vee (d_j-3)}.
\end{equation}
Finally, since for all $p\ge 1$, we have
\begin{equation*}
    p \left( \frac{d_\ell-1}{2} - 1 + \frac{d_\ell+d_j}{2} - 2 + \frac{d_j+3}{2} - d_j \right) \ge d_\ell-d_j \iff p \ge \frac{d_\ell-d_j}{d_\ell-2}.
\end{equation*}
Thus, by Lemma~\ref{Nix} again, $S_3$ is bounded in $L^p$ provided 
\begin{equation}\label{S3}
    \frac{2d_\ell}{2d_\ell \wedge (d_\ell+3)} < p < \frac{2d_j}{0 \vee (d_j-3)}.
\end{equation}
Note that $\left( \frac{2D}{D+1}, \frac{2D}{D-1} \right)$ is contained in each of the intervals in \eqref{S1}, \eqref{S2} and \eqref{S3}. The proof of Claim~1 is then complete.
\end{proof}


It remains to show that an operator with kernel $\frac{e^{\mp i(|x|+|y|)}}{|x|+|y|}$ can be bounded on $L^p$ only when $p=2$. That is, we have to establish the $L^p$-unboundedness of the linear combination of the following two operators:
\begin{equation*}
S_4f(x) := \sum_{\ell \ne j}|x|^{-\frac{d_\ell-1}{2}} \int_1^\infty |y|^{-\frac{d_j-1}{2}} \frac{e^{i(|x|+|y|)}f(y)}{|x|+|y|} |y|^{d_j-1} dy:= \sum_{\ell \ne j}S_4^{\ell j}f(x), 
\end{equation*}
and
\begin{equation*}
S_5f(x):= \sum_{\ell \ne j} |x|^{-\frac{d_\ell-1}{2}} \int_1^\infty |y|^{-\frac{d_j-1}{2}} \frac{e^{-i(|x|+|y|)}f(y)}{|x|+|y|} |y|^{d_j-1} dy:= \sum_{\ell \ne j}S_5^{\ell j}f(x)
\end{equation*}
for $p \ne 2$.

Without loss of generality, assume that \(d_1<d_2\).
Since the spectral projection is self-adjoint, it  suffices to
consider $p>2$.  A direct estimate shows  that the kernel of $S_4^{21}+S_{5}^{21}$
satisfies 
\begin{align*}
    |(S_4^{21}+S_{5}^{21})(x,y)| \lesssim \begin{cases}
        |x|^{-\frac{d_2-1}{2}} |y|^{-\frac{d_1+1}{2}}, & |y|\ge |x|,\\
        |x|^{-\frac{d_2+1}{2}} |y|^{-\frac{d_1-1}{2}}, & |y|\le |x|.
    \end{cases}
\end{align*}
Since 
\begin{equation*}
    p\left( \frac{d_2-1+d_1+1}{2} - d_1 \right) \ge d_2-d_1 \implies p\ge 2,
\end{equation*}
Lemma~\ref{Nix} tells us $S_4^{21}+S_{5}^{21}$ is bounded from $L^p([1,\infty), r^{d_1-1}dr) \to L^p([1,\infty), r^{d_2-1}dr)$ for 
\begin{equation*}
    2\le  p < \frac{2D}{D-1}.
\end{equation*}
Hence, the result of Theorem~\ref{fefferman_nondoubling} follows by showing the unboundedness of the linear combination of $S_4^{12}$ and $S_{5}^{12}$, say $c_1 S_4^{12} + c_2 S_5^{12}$ with $c_1,c_2 \in \mathbb C$ not simultaneously zero, from $L^p([1,\infty), r^{d_2-1}dr) \to L^p([1,\infty), r^{d_1-1}dr)$ for $p>2$ and $d_1<d_2$. 

Assume $c_1\ne 0$. Consider the function 
\begin{equation}\label{couterF}
    f_\varepsilon(y) = e^{-iy} \mathbf{1}_{[\varepsilon,2\varepsilon]}(y) y^{-\frac{d_2-1}{2}}, \quad \varepsilon\ge 1,
\end{equation}
and if $c_1=0$ (so $c_2\ne 0$), we put $f_\varepsilon(y) = e^{iy} \mathbf{1}_{[\varepsilon,2\varepsilon]}(y) y^{-\frac{d_2-1}{2}}$.

First,
\begin{align*}
    \| f_\varepsilon \|_p^p \lesssim \int_\varepsilon^{2\varepsilon} y^{-p \frac{d_2-1}{2}} y^{d_2-1} dy \lesssim \varepsilon^{d_2 - p \frac{d_2-1}{2}}.
\end{align*}
Second, for $x\in [2\varepsilon,3\varepsilon]$,
\begin{align*}
    |S_4^{12}f_{\varepsilon}(x)| = x^{-\frac{d_1-1}{2}} \int_\varepsilon^{2\varepsilon} \frac{dy}{x+y} \gtrsim \varepsilon^{-\frac{d_1-1}{2}}.
\end{align*}
Third, for $x\in [2\varepsilon,3\varepsilon]$, one also has
\begin{align*}
    |S_5^{12} f_\varepsilon(x)| &= x^{-\frac{d_1-1}{2}} \left| \int_\varepsilon^{2\varepsilon} \frac{e^{-2iy}dy}{x+y} \right| \\
    &= x^{-\frac{d_1-1}{2}} \left| -\frac{1}{2i} \int_\varepsilon^{2\varepsilon} \frac{d (e^{-2iy}) }{x+y} \right|\\
    &= \frac{1}{2} x^{-\frac{d_1-1}{2}} \left| \left[\frac{e^{-2iy}}{x+y}\right]_{y=\varepsilon}^{2\varepsilon} +   \int_\varepsilon^{2\varepsilon} \frac{e^{-2iy}}{(x+y)^2} dy  \right|\\
    &\lesssim \varepsilon^{-\frac{d_1+1}{2}}.
\end{align*} 
Hence, for sufficiently large $\varepsilon$ and  for $x\in [2\varepsilon,3\varepsilon]$ we have 
\begin{equation*}
    |c_1 S_4^{12}f_\varepsilon(x) + c_2 S_5^{12}f_\varepsilon(x)| \gtrsim |S_4^{12}f_\varepsilon(x)| \gtrsim \varepsilon^{-\frac{d_1-1}{2}}.
\end{equation*}
Thus,
\begin{align*}
 \| c_1 S_4^{12} f_\varepsilon + c_2S_5^{12} f_\varepsilon\|_{L^p([2\varepsilon,3\varepsilon], x^{d_1-1} dx)}^p \gtrsim \| c_1 S_4^{12} f_\varepsilon + c_2 S_5^{12} f_\varepsilon \|_{L^p([2\varepsilon,3\varepsilon], x^{d_1-1}dx)}^p \\ \gtrsim \int_{2\varepsilon}^{3\varepsilon} \varepsilon^{-p \frac{d_1-1}{2}} x^{d_1-1} dx \gtrsim \varepsilon^{d_1-p\frac{d_1-1}{2}}.
\end{align*}
Therefore, if $S_4+S_5$ is bounded in $L^p$ for some $p>2$, then we have for $\varepsilon  \gg 1$,
\begin{equation*}
    \varepsilon^{d_1-p\frac{d_1-1}{2}} \lesssim \varepsilon^{d_2 - p \frac{d_2-1}{2}} \iff \varepsilon^{(d_1-d_2)(1-p/2)} \lesssim 1,
\end{equation*}
which is not true as $\varepsilon\to \infty$. The proof of Theorem~\ref{fefferman_nondoubling} is now complete.
\end{proof}

\end{document}
